% TOP_PROBLEM: {"id": 1453, "title": "Finite Lattice Representation Problem", "category_id": 4, "category": {"id": 4, "name": "algebra", "display_name": "Algebra", "description": "Group theory, ring theory, field theory, and algebraic structures.", "slug": "algebra", "order_index": 4, "created_at": "2026-07-31T15:26:25.671Z"}, "status": "open"}
% FIRST_POSTED: 2026-10-01
% ATTEMPT_STATUS: unresolved
% Imported from: unsolvedmath_additional_500.tex; UnsolvedMath ID 1453
% !TeX program = lualatex
% Compile twice: lualatex 1453.tex
% Self-contained: no external bibliography, images, or \input files.
% Numeric section labels and visible numbers are original UnsolvedMath IDs.
\documentclass[11pt,a4paper]{article}
\usepackage[margin=24mm,headheight=14pt]{geometry}
\usepackage{fontspec}
\setmainfont{Latin Modern Roman}
\setsansfont{Latin Modern Sans}
\setmonofont{Latin Modern Mono}
\usepackage{amsmath,amssymb,mathtools}
\usepackage{unicode-math}
\setmathfont{Latin Modern Math}
\usepackage{microtype}
\usepackage{enumitem,xurl,needspace}
\usepackage[dvipsnames]{xcolor}
\usepackage{hyperref}
\hypersetup{unicode=true,colorlinks=true,linkcolor=MidnightBlue,urlcolor=MidnightBlue,
 pdftitle={UnsolvedMath Problem 1453},pdfauthor={Source-annotated compilation},bookmarksnumbered=true}
\usepackage{bookmark,tocloft,titlesec,fancyhdr}
\setlength{\cftsecnumwidth}{4.2em}
\cftsetpnumwidth{2.5em}
\cftsetrmarg{4em}
\titleformat{\section}{\Large\bfseries\raggedright}{\thesection}{0.7em}{}
\pagestyle{fancy}\fancyhf{}
\fancyhead[L]{\small\sffamily Additional UnsolvedMath statements}
\fancyhead[R]{\small Original catalogue IDs}
\fancyfoot[C]{\thepage}
\setlength{\parindent}{0pt}
\setlength{\parskip}{5pt plus 1pt}
\setlength{\emergencystretch}{3em}
\setcounter{tocdepth}{1}\setcounter{secnumdepth}{1}
\setlist[itemize]{leftmargin=3em,itemsep=4pt,topsep=4pt,parsep=0pt}
\allowdisplaybreaks
\newcommand{\N}{\mathbb{N}}
\newcommand{\Z}{\mathbb{Z}}
\newcommand{\Q}{\mathbb{Q}}
\newcommand{\R}{\mathbb{R}}
\newcommand{\C}{\mathbb{C}}
\newcommand{\F}{\mathbb{F}}
\newcommand{\A}{\mathbb{A}}
\newcommand{\PP}{\mathbb{P}}
\newcommand{\E}{\mathbb{E}}
\newcommand{\eps}{\varepsilon}
\newcommand{\dd}{\,\mathrm d}
\newcommand{\sref}[2]{\hyperlink{source:#1:#2}{[S#2]}}
\newif\ifshowcatalogue
\showcataloguetrue
\newcommand{\cataloguescope}[1]{\ifshowcatalogue
 \paragraph{Statement recorded in the supplied source.}
 #1\fi}
\begin{document}
% UnsolvedMath ID 1453; source catalogue code ALG-008

\setcounter{section}{1452}

\section{Finite Lattices as Congruence Lattices}\label{1453}

{\small\sffamily UnsolvedMath ID 1453 · Algebra}

\subsection{Definitions and mathematical statement}

\paragraph{Shared notation.}Unless a section says otherwise, $\N=\{1,2,\ldots\}$,
$\Z$ is the ring of integers, $\Q,\R,\C$ are the rational, real, and complex
fields, $[N]=\{1,\ldots,\lfloor N\rfloor\}$, and $\log$ is the natural
logarithm. For real functions of a parameter tending to infinity,
$f\ll g$ and $f=O(g)$ mean $|f|\leq Cg$ eventually for some fixed $C$;
$f\gg g$ reverses this comparison; $f=o(g)$ means $f/g\to0$;
$f\sim g$ means $f/g\to1$; and $f\asymp g$ means both order bounds.
A subscript on $O$ or $\ll$ specifies allowed constant dependence.
The expression $n^{o(1)}$ in an upper bound means growth smaller than
$n^\epsilon$ for each fixed $\epsilon>0$. Local definitions govern any
exception, finite-field convention, or use of the same symbol for another
object. An asymptotic claim is not automatically an effective algorithm.



A finite algebra is a finite nonempty set equipped with finitary operations. A congruence is an equivalence relation preserved by every operation. The set $\operatorname{Con}(A)$ of congruences forms a lattice under inclusion, with meet given by intersection. Ask whether each finite lattice is isomorphic, as a lattice, to $\operatorname{Con}(A)$ for some finite algebra. \sref{1453}{1} \sref{1453}{2}

\paragraph{Statement recorded in the supplied source.}
Is every finite lattice isomorphic to the congruence lattice of some finite algebra? \sref{1453}{1}

\paragraph{Residual task reported by the archive.}

The embedded review (2026-08-17) records: Represent every finite lattice or prove a finite nonrepresentable lattice. \sref{1453}{1}

\subsection{Short English statement}

Can every finite lattice arise as the collection of operation-compatible equivalence relations on a finite algebra?

\subsection{Research attempt: unary reduction, finite certificates, and explicit lattice families}
\textbf{Status.} We give a finite-search formulation and several infinite families of representable lattices, not a representation of every finite lattice. The primary paper [R1] studies small unary representations; the reduction below is established directly rather than inferred from its abstract.

\subsubsection{1. Arbitrary operations can be replaced by unary translations}
Let $A$ be a finite algebra. For each basic operation $f$ and each input coordinate, fix arbitrary elements in all other coordinates. This produces a unary map $u:A\to A$, called a translation. An equivalence relation preserved by every operation is preserved by all these translations. Conversely, suppose an equivalence relation $\theta$ is preserved by every translation. Given tuples $(a_1,\ldots,a_n)$ and $(b_1,\ldots,b_n)$ with $a_i\mathrel\theta b_i$, change their coordinates one at a time. Preservation by the corresponding translation shows the successive outputs are $\theta$-equivalent, and transitivity gives $f(a_1,\ldots,a_n)\mathrel\theta f(b_1,\ldots,b_n)$. Constants impose no extra equivalence conditions. Thus the original congruence lattice equals that of its unary translation algebra.

If $|A|=m$, there are only $m^m$ unary functions in total. Therefore every representable lattice has a representation on the same finite carrier using a finite subset of these functions, even if its initially supplied signature was unrestricted.

\subsubsection{2. An exhaustive positive search, with a precise stopping limitation}
For $m=1,2,\ldots$, enumerate all subsets $U$ of the $m^m$ unary maps and all partitions of an $m$-element set. A partition $\theta$ is a congruence exactly when
\[a\mathrel\theta b\Longrightarrow u(a)\mathrel\theta u(b)
\quad\text{for every }u\in U\text{ and every }a,b.
\]
This is a finite test. Order the retained partitions by refinement and test isomorphism of this finite ordered set with the input lattice. Order isomorphism suffices because meets and joins are determined by the order. Any successful representation is an explicit finite certificate. The unary reduction proves that the search eventually succeeds whenever any finite algebra represents the lattice.

It does not provide a carrier bound in terms of the input lattice, so a failed search up to some $m$ is not a nonrepresentability proof. Enumerability of the positive cases is weaker than a solution to the conjecture.

\subsubsection{3. Groups give concrete families}
For a group viewed as an algebra with multiplication, inverse and identity, a congruence is determined by its identity class $N$, which is a normal subgroup. Indeed $a\mathrel\theta b$ exactly when $ab^{-1}\in N$, and normality follows from compatibility with conjugation. Conversely each normal subgroup defines such a congruence.

For $C_n$, the normal-subgroup lattice is the divisor lattice of $n$. Writing $n=\prod_{i=1}^r p_i^{a_i}$ identifies it with the product of chains $\prod_i\{0,\ldots,a_i\}$. Thus every finite chain and every finite product of finite chains occurs; squarefree $n$ gives Boolean lattices. For the additive group $\mathbb F_p^2$, its subgroups are vector subspaces. It has $p+1$ one-dimensional subspaces, since each contains $p-1$ nonzero vectors and these partition the $p^2-1$ nonzero vectors. Their lattice consists of a bottom, a top, and $p+1$ incomparable intermediate points, giving another explicit family.

\subsubsection{4. Verification and missing universality}
The diagnostic checks unary preservation against binary-operation preservation for every binary operation on a two-element carrier and every equivalence relation there. It also counts the one-dimensional subspaces of $\mathbb F_p^2$ for $p=2,3,5,7$. The group constructions do not cover every finite lattice, and the unary reduction changes the operation set without enlarging the carrier only after an algebra already exists. Using that reduction as an existence proof would reverse its logical direction. A universal carrier construction or an obstruction valid for all carrier sizes remains absent.

\subsection{Sources}

{\small

\begin{itemize}

\item[\hypertarget{source:1453:1}{S1}] ULAM.AI, \emph{UnsolvedMath}, supplied \texttt{problems.json}, record 1453 (ALG-008), statement and background. The embedded literature-review date is 2026-08-17; that is the archive’s review date, not a fresh certification. Dataset landing page: \url{https://huggingface.co/datasets/ulamai/UnsolvedMath}.

\item[\hypertarget{source:1453:2}{S2}] Finite lattice representation problem overview (). \url{https://en.wikipedia.org/wiki/Finite_lattice_representation_problem}. Bibliographic information imported from the supplied record.

\item[R1] R. Bunn, D. Grow, M. Insall and P. Thiem, A minimal congruence lattice representation for $M_{p+1}$, Journal of the Australian Mathematical Society 108 (2020), 332--340.
\url{https://www.cambridge.org/core/journals/journal-of-the-australian-mathematical-society/article/minimal-congruence-lattice-representation-for-mathbbmp1/F1031B09770B7060128554C0A2E86EB5}.
\end{itemize}

}

\label{end:1453}
\end{document}
